%%--------------------------------------------------------------------
%1--------------------------------------------------------------------
\begin{glyph}{Attach (付)}{attach}\label{glyph:attach}
Attach, in both physical and figurative senses.\\\hline
This 漢字 conveys the meanings of many English words---from \textit{setting} something on fire, to \textit{turning on} an electrical device.
\end{glyph}
\gindex{Attach}{付}\strindex{5}{付}

\begin{comps}
\comprtwocone&亻 + 寸\\\hline
\comprthreecone&Person + Tiny/Inch\\\hline
\end{comps}

\begin{remglyph}
\hooglig{Attach} the \remcom{human} to the \remcom{tiny} parasite.\\\hline
\end{remglyph}

\begin{prons}
\pronsrtwocone&\textbf{フ (FU)}&\textbf{つ (tsu)}\\\hline
\rye{3}{The reading in the second position is つ when 付 is preceded by another \mbox{漢字} that is read with its verb stem: \mbox{見付ける} (み{\wsplit}つ{\middel}ける: \textit{to find}) is one example.\\\hline
When the preceding reading does not form a verb with its 漢字, 付 will sometimes becomes voiced as づ: e.g. \mbox{気付く} (キ{\wsplit}づ{\middel}く: \textit{to notice}) as 気 does not form a standalone verb with キ.}
\pronsrthreecone&---&---\\\hline
\end{prons}

\begin{remprons}
Getting \hooglig{attached} to \comon{foosball} can stir up a \comkun{tsunami}.\\\hline
\end{remprons}

%{\middel}
% &  ()\\\hline
\begin{somme}
\splitter{付く}{intr}&\textit{to be attached}&つ{\middel}く \mbox{(tsu{\middel}ku)} \\\hline
\splitter{付ける}{tr}&\textit{to attach/put on}&つ{\middel}ける \mbox{(tsu{\middel}keru)} \\\hline
\splitter{見付ける}{tr}&see + attach = \textit{to find}&み{\wsplit}つける \mbox{(mi{\wsplit}tsu{\middel}keru)} \\\hline
\splitter{見せ付ける}{tr}&see + attach = \textit{to make a display of}&み{\middel}せ{\wsplit}つ{\middel}ける \mbox{(mi{\middel}se{\wsplit}tsu{\middel}keru)}\\\hline
\splitter{気付く}{tr}&spirit + attach = \textit{to notice (sth.)}&キ{\wsplit}づ{\middel}く \mbox{(KI{\wsplit}zu{\middel}ku)}\\\hline
\splitter{名付ける}{tr}&name + attach = \textit{to name}&な{\wsplit}づ{\middel}ける \mbox{(na{\wsplit}zu{\middel}keru)} \\\hline
\splitter{近付ける}{intr}&near + attach = \textit{to bring near}&ちか{\wsplit}づ{\middel}ける \mbox{(chika{\wsplit}zu{\middel}keru)}\\\hline
\splitter{受け付ける}{tr}&receive + attach = \textit{to accept/receive}&う{\middel}け{\wsplit}つ{\middel}ける \mbox{(u{\middel}ke{\wsplit}tsu{\middel}keru)}\\\hline
\splitter{言付ける}{tr}&say + attach = \textit{to order/direct}&いい{\wsplit}つ{\middel}ける \mbox{(\={i}{\wsplit}tsu{\middel}keru)}\\\hline
\splitter{引き付ける}{tr}&pull + attach = \textit{to draw near}&ひ{\middel}き{\wsplit}つ{\middel}ける \mbox{(hi{\middel}ki{\wsplit}tsu{\middel}keru)}\\\hline
\splitter{押し付ける}{tr}&push + attach = \textit{to press against}&お{\middel}し{\wsplit}つ{\middel}ける \mbox{(o{\middel}shi{\wsplit}tsu{\middel}keru)}\\\hline
\splitter{押さえ付ける}{tr}&push + attach = \textit{to hold down}&お{\middel}さえ{\wsplit}つ{\middel}ける \mbox{(o{\middel}sae{\wsplit}tsu{\middel}keru)}\\\hline
\splitter{売り付ける}{tr}&sell + attach = \textit{to force a sale}&う{\middel}り{\wsplit}つ{\middel}ける \mbox{(u{\middel}ri{\wsplit}tsu{\middel}keru)}\\\hline
\splitter{聞き付ける}{tr}&hear + attach = \textit{to hear; to catch the sound of}&き{\middel}き{\wsplit}つ{\middel}ける \mbox{(ki{\middel}ki{\wsplit}tsu{\middel}keru)}\\\hline
付近&attach + near = \textit{neighbourhood}&フ{\middel}キン \mbox{(FU{\middel}KIN)} \\\hline
受付&receive + attach = \textit{receipt; reception desk}&うけ{\middel}つけ \mbox{(uke{\middel}tsuke)} \\\hline
\splitterny{貸し付ける}{tr}&lend + attach = \textit{to lend/loan}&か{\middel}し{\wsplit}つ{\middel}ける \mbox{(ka{\middel}shi{\wsplit}tsu{\middel}keru)}\\\hline
\splitterny{結び付ける}{tr}&bind + attach = \textit{to combine}&むす{\middel}び{\wsplit}つ{\middel}ける \mbox{(musu{\middel}bi{\wsplit}tsu{\middel}keru)}\\\hline
\end{somme}
\end{multicols}

\begin{decomp}{6}
これ&を&受付&に&お出し&下さい。\\\hline
これ&を&うけつけ&に&おだし&ください\\\hline
this&を&reception desk&に&submit&please\\\hline
\multicolumn{6}{|r|}{\textit{Please hand this in at the front desk.}}\\\hline
\end{decomp}
\pagebreak

%%--------------------------------------------------------------------
%2--------------------------------------------------------------------

\begin{multicols}{2}
\begin{glyph}{Mourn (弔)}{mourn}\label{glyph:mourn}
Mourn.
\end{glyph}
\gindex{Mourn}{弔}\strindex{4}{弔}

\begin{comps}
\comprtwocone&弓 + 丨\\\hline
\comprthreecone&Bow + Pole\\\hline
\end{comps}

\begin{remglyph}
\hooglig{Mourn} the \remcom{bow} that turned into a \remcom{pole}.\\\hline
\end{remglyph}

\begin{prons}
\pronsrtwocone&\textbf{チョウ (CH\={O})}&---\\\hline
\pronsrthreecone&--- &とうら (tomura)\\\hline
\end{prons}

\begin{remprons}
I \hooglig{mourned} about the \comon{chores} such that \ucomkun{tomme luxury} could not cheer me up.\\\hline
\end{remprons}

%{\middel}
% &  ()\\\hline
\begin{somme}
弔意& mourn + mind = \textit{condolence}&チョウ{\middel}イ \mbox{(CH\={O}{\middel}I)}\\\hline
弔文& mourn + literature = \textit{funeral address}&チョウ{\middel}ブン \mbox{(CH\={O}{\middel}BUN)}\\\hline
\notyet{弔問}& mourn + question = \textit{sympathy call}&チョウ{\middel}モン \mbox{(CH\={O}{\middel}MON)}\\\hline
\end{somme}

\begin{decomp}{4}
弔い&は&今日&です。\\\hline
とむらい&は&きょう&です\\\hline
funeral&は&today&be/is\\\hline
\multicolumn{4}{|r|}{\textit{The funeral is today.}}\\\hline
\end{decomp}
\end{multicols}

%%--------------------------------------------------------------------
%3--------------------------------------------------------------------
\begin{multicols}{2}
\begin{glyph}{Cry (泣)}{cry}\label{glyph:cry}
Cry.
\end{glyph}
\gindex{Cry}{泣}\strindex{8}{泣}

\begin{comps}
\comprtwocone&氵 + 立\\\hline
\comprthreecone&Water + Stand\\\hline
\end{comps}

\begin{remglyph}
\hooglig{Cry} \remcom{water} as I \remcom{stand}.\\\hline
\end{remglyph}

\begin{prons}
\pronsrtwocone&---&\textbf{な (na)}\\\hline
\pronsrthreecone&キュウ (KY\={U})&---\\\hline
\end{prons}

\begin{remprons}
\hooglig{Crying} myself \comkun{numb} from all the \ucomon{`kill you'} threats.\\\hline
\end{remprons}

%{\middel}
% &  ()\\\hline
\begin{somme}
泣く&\textit{cry}&な{\middel}く \mbox{(na{\middel}ku)}\\\hline
泣き出す&cry + exit = \textit{burst out crying}&な{\middel}き{\wsplit}だ{\middel}す \mbox{(na{\middel}ki{\wsplit}da{\middel}su)}\\\hline
泣き声&cry + voice = \textit{tearful voice}&な{\middel}き{\wsplit}ごえ \mbox{(na{\middel}ki{\wsplit}goe)}\\\hline
\notyet{泣き虫}&cry + insect = \textit{crybaby}&な{\middel}き{\wsplit}むし \mbox{(na{\middel}ki{\wsplit}mushi)}\\\hline
\end{somme}

\begin{decomp}{3}
母&は&泣いていた。\\\hline
はは&は&ないていた\\\hline
mother&は&cry\\\hline
\multicolumn{3}{|r|}{\textit{My mother was crying.}}\\\hline
\end{decomp}
\end{multicols}

\pagebreak
%%--------------------------------------------------------------------
%4--------------------------------------------------------------------

\begin{multicols}{2}
\begin{glyph}{Silver (銀)}{silver}\label{glyph:silver}
Silver.\\\hline
Recall from Entry \ref{glyph:gold} (page \pageref{glyph:gold}) that when the character for gold functions as a component in other 漢字, it means `metal' in general.
\end{glyph}
\gindex{Silver}{銀}\strindex{14}{銀}

\begin{comps}
\comprtwocone & 金 + 艮\\\hline
\comprthreecone & Metal + Boundary\\\hline
\end{comps}

\begin{remglyph}
\hooglig{Silver} \remcom{metal} \remcom{boundaries}.\\\hline
\end{remglyph}

\begin{prons}
\pronsrtwocone & \textbf{ギン (GIN)} & --- \\\hline
\pronsrthreecone & --- & --- \\\hline
\end{prons}

\begin{remprons}
\hooglig{Silver} \comon{ginkgo} capsules.\\\hline
{\warn}\textit{Ginkgo}, also known as maidenhair tree, is a species of tree native to China.\\\hline
\end{remprons}

%{\middel}
% &  ()\\\hline
\begin{somme}
銀  & \textit{silver}  & ギン \mbox{(GIN)}  \\\hline
銀行  & silver + go = \textit{bank}  & ギン{\middel}コウ \mbox{(GIN{\middel}K\={O})}  \\\hline
\notyet{銀世界}  & silver + society + world = \textit{snow-covered scene}  & ギン{\middel}セ{\middel}カイ \mbox{(GIN{\middel}SE{\middel}KAI)}\\\hline
\end{somme}

\begin{decomp}{5}
金&は&銀&より&重い。\\\hline
キン&は&ギン&より&おもい\\\hline
gold&は&silver&more&heavy\\\hline
\multicolumn{5}{|r|}{\textit{Gold is heavier than silver.}}\\\hline
\end{decomp}
\end{multicols}

%%--------------------------------------------------------------------
%5--------------------------------------------------------------------
\begin{multicols}{2}
\begin{glyph}{Question (問)}{question}\label{glyph:question}
Question.  Inquiry.\\\hline
In some words (the second compound below is an example), this 漢字 can convey the sense of \textit{problem}.
\end{glyph}
\gindex{Question}{問}\strindex{11}{問}

\begin{comps}
\comprtwocone & 門 + 口 \\\hline
\comprthreecone & Gate + Mouth \\\hline
\end{comps}

\begin{remglyph}
\hooglig{Questions} have their \remcom{gateway} through the \remcom{mouth}.\\\hline
\end{remglyph}

\begin{prons}
\pronsrtwocone & \textbf{モン (MON)} & \textbf{と (to)}\\\hline
\pronsrthreecone & --- & --- \\\hline
\end{prons}

\begin{remprons}
\hooglig{Questionable} \comkun{torrents} on the \comon{monastery}.\\\hline
\end{remprons}

%{\middel}
% &  ()\\\hline
\begin{somme}
質問  & quality + question = \textit{question}  & シツ{\middel}モン \mbox{(SHITSU{\middel}MON)}  \\\hline
\notyet{問題}  & question + topic = \textit{problem; issue}  & モン{\middel}ダイ \mbox{(MON{\middel}DAI)}  \\\hline
\notyet{問い合わせる}  & question + match = \textit{inquire; apply to}  & と{\middel}い　あ{\middel}わせる \mbox{(to{\middel}i a{\middel}waseru)}  \\\hline
\end{somme}

\begin{irrr}
\notyet{問屋}  & question + roof = \textit{wholesale dealer}  & とん{\middel}や (ton{\middel}ya)  \\\hline
\end{irrr}

\begin{decomp}{5}
其れ&は&問題外&だ&よ。\\\hline
それ&は&モンダイガイ&だ&よ\\\hline
that&は&unthinkable&be/is&よ\\\hline
\multicolumn{5}{|r|}{\textit{That is out of the question.}}\\\hline
\end{decomp}
\end{multicols}

\pagebreak
%%--------------------------------------------------------------------
%6--------------------------------------------------------------------

\begin{multicols}{2}
\begin{glyph}{Dark (暗)}{dark}\label{glyph:dark}
Dark, with its various implications.
\end{glyph}
\gindex{Dark}{暗}\strindex{13}{暗}

\begin{comps}
\comprtwocone & 日 + 音 \\\hline
\comprthreecone & Sun/Day + Sound \\\hline
\end{comps}

\begin{remglyph}
\hooglig{Dark} and ominous \remcom{sun} \remcom{sounds}.\\\hline
\end{remglyph}

\begin{prons}
\pronsrtwocone & \textbf{アン (AN)} & \textbf{くら (kura)}\\\hline
\pronsrthreecone & --- & --- \\\hline
\end{prons}

\begin{remprons}
\hooglig{Dark} \comkun{curator} of an \comon{uncle}.\\\hline
\end{remprons}

%{\middel}
% &  ()\\\hline
\begin{somme}
暗い  & \textit{dark}  & くら{\middel}い \mbox{(kura{\middel}i)}  \\\hline
暗黒  & dark + black = \textit{darkness; blackness}  & アン{\middel}コク \mbox{(AN{\middel}KOKU)}  \\\hline
\notyet{暗記}  & dark + note = \textit{memorization}  & アン{\middel}キ \mbox{(AN{\middel}KI)}  \\\hline
\end{somme}

\begin{decomp}{4}
外&が&暗く&なってきた。\\\hline
そと&が&くらく&なってきた\\\hline
outside&が&dark&to grow\\\hline
\multicolumn{4}{|r|}{\textit{It's getting darker outside now.}}\\\hline
\end{decomp}
\end{multicols}

%%--------------------------------------------------------------------
%7--------------------------------------------------------------------
\begin{multicols}{2}
\begin{glyph}{Picture (絵)}{picture}\label{glyph:picture}
Picture.
\end{glyph}
\gindex{Picture}{絵}\strindex{12}{絵}

\begin{comps}
\comprtwocone & 糸 + 会 \\\hline
\comprthreecone & Thread + Meet (\ref{glyph:meet}) \\\hline
\end{comps}

\begin{remglyph}
\hooglig{Picturesque} \remcom{threads} \remcom{meet}.\\\hline
\hooglig{Picture} a place where \remcom{threads} \remcom{meet} artistically.\\\hline
\end{remglyph}

\begin{prons}
\pronsrtwocone & \textbf{エ (E)} & --- \\\hline
\pronsrthreecone & カイ (KAI) & --- \\\hline
\rye{3}{This reading appears in only one compound.  The word, however, is commonly used: 絵画 (カイ{\middel}ガ) \textit{pictures and paintings}, composed of, naturally enough, the 漢字 for `picture' and `painting' (Entry \ref{glyph:painting}, page \pageref{glyph:painting}).}
\end{prons}

\begin{remprons}
\hooglig{Picture} an \comon{egg} \ucomon{kite}.\\\hline
\end{remprons}

%{\middel}
% &  ()\\\hline
\begin{somme}
絵 & \textit{picture}  & エ \mbox{(E)}\\\hline
絵本 & picture + main (book) = \textit{picture book}  & エ{\middel}ホン \mbox{(E{\middel}HON)}\\\hline
絵画 & picture + painting = \textit{pictures and paintings} & カイ{\middel}ガ　\mbox{(KAI{\middel}GA)}\\\hline
\end{somme}

\begin{decomp}{5}
壁&の&絵&を&見て。\\\hline
かべ&の&エ&を&みて\\\hline
wall&の&painting&を&look at\\\hline
\multicolumn{5}{|r|}{\textit{Look at the picture on the wall.}}\\\hline
\end{decomp}
\end{multicols}

\pagebreak
%%--------------------------------------------------------------------
%8--------------------------------------------------------------------

\begin{multicols}{2}
\begin{glyph}{Surname (氏)}{surname}\label{glyph:surname}
Surname.\\\hline
This 漢字 can also mean \textit{Clan} in some instances, as well as convey a more official-sounding form of \textit{Mr.} when it is attached to a family name (as in the final compound below).
\end{glyph}
\gindex{Surname}{氏}\strindex{4}{氏}

\begin{comps}
\comprtwocone &  氏\\\hline
\comprthreecone & Surname/Clan\\\hline
\end{comps}

\begin{remglyph}
 Part of the 漢字 radicals (\#{83}).\\\hline
\end{remglyph}

\begin{prons}
\pronsrtwocone & \textbf{シ (SHI)} & --- \\\hline
\pronsrthreecone & --- & うじ (uji) \\\hline
\end{prons}

\begin{remprons}
I \hooglig{surnamed} my \comon{shield} according to the \ucomkun{uji} customs.\\\hline
{\warn}\textit{Uji} is a city in Ky\={o}to.\\\hline
\end{remprons}

%{\middel}
% &  ()\\\hline
\begin{somme}
氏名  & surname + name = \textit{one's full name}  & シ{\middel}メイ \mbox{(SHI{\middel}MEI)}  \\\hline
中村氏  & middle + village + surname = \textit{Mr. Nakamura}  & なか{\middel}むら{\middel}シ \mbox{(naka{\middel}mura{\middel}SHI)}  \\\hline
\end{somme}

\begin{decomp}{7}
あなた&の&氏名&は&何&です&か\\\hline
あなた&の&シメイ&は&なに&です&か\\\hline
you&の&full name&は&what&be/is&か\\\hline
\multicolumn{7}{|r|}{\textit{What's your name?}}\\\hline
\end{decomp}
\end{multicols}

%%--------------------------------------------------------------------
%9-------------------------------------------------------------------
\begin{multicols}{2}
\begin{glyph}{Young (若)}{young}\label{glyph:young}
Young.
\end{glyph}
\gindex{Young}{若}\strindex{8}{若}

\begin{comps}
\comprtwocone & 艹 + 右 \\\hline
\comprthreecone & Grass + Right\\\hline
\end{comps}

\begin{remglyph}
\hooglig{Young} \remcom{grass} on the \remcom{right}-hand side.\\\hline
\end{remglyph}

\begin{prons}
\pronsrtwocone & --- & \textbf{わか (waka)}\\\hline
\pronsrthreecone & ジャク、ニャク (JAKU, NYAKU) & --- \\\hline
\end{prons}

\begin{remprons}
\hooglig{Young} \comkun{wakaberry} enthusiasts \ucomon{jucked} away his cup from the \ucomon{`knee-yuck'} flavour.\\\hline
\end{remprons}

%{\middel}
% &  ()\\\hline
\begin{somme}
\rye{3}{Note that the second compound uses the less common pronunciation for 者.}
若い  & \textit{young}  & わか{\middel}い \mbox{(waka{\middel}i)}  \\\hline
若者  & young + individual = \textit{young person}  & わか{\middel}もの \mbox{(waka{\middel}mono)}  \\\hline
若々しい  & young + young = \textit{youthful}  &  わか{\middel}わか{\middel}しい \mbox{(waka{\middel}waka{\middel}shii)} \\\hline
\end{somme}

\begin{irrr}
若人  & young + person = \textit{young person}  & わこう{\middel}ど \mbox{(wak\={o}{\middel}do)} \\\hline
\end{irrr}

\begin{decomp}{3}
僕&は&若い。\\\hline
ボク&は&わかい\\\hline
I&は&young\\\hline
\multicolumn{3}{|r|}{\textit{I'm young.}}\\\hline
\end{decomp}
\end{multicols}

\begin{decomp}{7}
此れ&が&若者&特有&の&欠点&です。\\\hline
これ&が&わかもの&トクユウ&の&ケッテン&です\\\hline
this&が&youth&peculiar to&の&defect&be/is\\\hline
\multicolumn{7}{|r|}{\textit{This is a weakness peculiar to young people.}}\\\hline
\end{decomp}

\pagebreak
%%--------------------------------------------------------------------
%10-------------------------------------------------------------------

\begin{multicols}{2}
\begin{glyph}{Cause (由)}{cause}\label{glyph:cause}
Cause.  Reason.\\\hline
This is written in the same order as the 漢字 for ``Rice Field'' (Entry \ref{glyph:rice_field}, page \pageref{glyph:rice_field}), but with the third stroke lengthened.
\end{glyph}
\gindex{Cause}{由}\strindex{5}{由}

\begin{comps}
\comprtwocone & 田 \\\hline
\comprthreecone & Rice Field \\\hline
\end{comps}

\begin{remglyph}
What \hooglig{caused} the \remcom{sun} to get \remcom{polarized}?\\\hline
\end{remglyph}

\begin{prons}
\pronsrtwocone & \textbf{ユ、ユウ (YU, Y\={U})} & --- \\\hline
\pronsrthreecone & ユイ (YUI) & よし (yoshi) \\\hline
\end{prons}

\begin{remprons}
\hooglig{`Cause} \ucomon{you insulted} the \comon{youthful} \ucomkun{yoshi} on \comon{Youtube}.\\\hline
{\warn}\textit{Yoshi} is an アニメ　character.\\\hline
\end{remprons}

%{\middel}
% &  ()\\\hline
\begin{somme}
\rye{3}{This character only appears in about 10 Japanese compounds, which makes the large number of readings for it seem extreme.  The examples below, however, are all very common.}
由来  & cause + come = \textit{origin}  & ユ{\middel}ライ \mbox{(YU{\middel}RAI)}\\\hline
自由  & self + cause = \textit{freedom}  & ジ{\middel}ユウ \mbox{(JI{\middel}Y\={U})}\\\hline
理由  & reason + cause = \textit{reason}  & リ{\middel}ユウ \mbox{(RI{\middel}Y\={U})}\\\hline
\end{somme}
\end{multicols}

\begin{decomp}{7}
多く&の&英単語&は&ラテン語&に&由来する。\\\hline
おおく&の&エイタンゴ&は&ラテンゴ&に&ユライする\\\hline
many&の&English words&は&Latin&に&origin\\\hline
\multicolumn{7}{|r|}{\textit{Many English words derive from Latin.}}\\\hline
\end{decomp}

%%--------------------------------------------------------------------
%11-------------------------------------------------------------------
\begin{multicols}{2}
\begin{glyph}{Declare (申)}{declare}\label{glyph:declare}
\textit{Declare}, in the sense of a more humble way to say the verb \textit{say}.\\\hline
Ensure that you are able to differentiate between this \mbox{漢字} and the preceding
entry for 由, as well as between these two and the character for ``Shell''
(甲), from Entry \ref{glyph:shell} (page \pageref{glyph:shell}).
\end{glyph}
\gindex{Declare}{申}\strindex{5}{申}

\begin{comps}
\comprtwocone & 田 \\\hline
\comprthreecone & Rice Field \\\hline
\end{comps}

\begin{remglyph}
\hooglig{Declare} the \remcom{sun's} sovereignty as \remcom{poles} of sunlight impale your skin.\\\hline
\end{remglyph}

\begin{prons}
\pronsrtwocone & \textbf{シン (SHIN)} & \textbf{もう (m\={o})}\\\hline
\pronsrthreecone & --- & --- \\\hline
\end{prons}

\begin{remprons}
\hooglig{Declare} \comon{sheen} \comkun{morphing}.\\\hline
\end{remprons}

%{\middel}
% &  ()\\\hline
\begin{somme}
申す& \textit{say (humble form)}  & もう{\middel}す \mbox{(m\={o}{\middel}su)}\\\hline
申し上げる& declare + upper = \textit{inform; tell (humble form)}  & もう{\middel}し　あ{\middel}げる \mbox{(m\={o}{\middel}shi a{\middel}geru)}\\\hline
内申& inside + declare = \textit{confidential report}  & ナイ{\middel}シン \mbox{(NAI{\middel}SHIN)}\\\hline
\end{somme}

\begin{decomp}{5}
私&は&周&と&申します。\\\hline
わたし&は&シュウ&と&もうします\\\hline
I&は&Sh\={u}&と&have the honour to\\\hline
\multicolumn{5}{|r|}{\textit{My surname is Zhou.}}\\\hline
\end{decomp}
\end{multicols}

\pagebreak
%%--------------------------------------------------------------------
%12-------------------------------------------------------------------

\begin{multicols}{2}
\begin{glyph}{Essential (要)}{essential}\label{glyph:essential}
Essential.  Necessary.
\end{glyph}
\gindex{Essential}{要}\strindex{9}{要}

\begin{comps}
\comprtwocone & 襾覀 + 女\\\hline
\comprthreecone & West + Woman\\\hline
\end{comps}

\begin{remglyph}
\hooglig{Essential} \remcom{western} products for \remcom{women}.\\\hline
\end{remglyph}

\begin{prons}
\pronsrtwocone & \textbf{ヨウ (Y\={O})} & \textbf{い (i)}\\\hline
\pronsrthreecone & --- & --- \\\hline
\end{prons}

\begin{remprons}
\hooglig{Essentially} an \comkun{interesting} \comon{yawn}.\\\hline
\end{remprons}

%{\middel}
% &  ()\\\hline
\begin{somme}
要る  & \textit{need; be needed}  & い{\middel}る \mbox{(i{\middel}ru)}  \\\hline
{\diff}要する  & \textit{require; need}  & ヨウ{\middel}する \mbox{(Y\={O}{\middel}suru)} \\\hline
\notyet{必要}  & certain + essential = \textit{essential; necessary}  & ヒツ{\middel}ヨウ \mbox{(HITSU{\middel}Y\={O})}  \\\hline
\notyet{重要}  & heavy + essential = \textit{importance}  & ジュウ{\middel}ヨウ \mbox{(J\={U}{\middel}Y\={O})}  \\\hline
\end{somme}

\begin{decomp}{7}
その&家&は&大&修理&を&要する。\\\hline
その&いえ&は&ダイ&シュウリ&を&ヨウする\\\hline
that&house&は&big&repairs&を&require\\\hline
\multicolumn{7}{|r|}{\textit{This house requires repairs.}}\\\hline
\end{decomp}

\begin{decomp}{5}
もっと&時間&が&必要&だ。\\\hline
もっと&ジカン&が&ヒツヨウ&だ\\\hline
more&time&が&necessary&be/is\\\hline
\multicolumn{5}{|r|}{\textit{I need more time}}\\\hline
\end{decomp}
\end{multicols}

%%--------------------------------------------------------------------
%13-------------------------------------------------------------------
\begin{multicols}{2}
\begin{glyph}{Reach (至)}{reach}\label{glyph:reach}
\textit{Reach}, in the sense of \textit{arrive at} or \textit{extend to}.
\end{glyph}
\gindex{Reach}{至}\strindex{6}{至}

\begin{comps}
\comprtwocone & 至\\\hline
\comprthreecone & Arrive\\\hline
\end{comps}

\begin{remglyph}
Part of the 漢字 radicals (\#{133}).\\\hline
\end{remglyph}

\begin{prons}
\pronsrtwocone & \textbf{シ (SHI)} & \textbf{いた (ita)}\\\hline
\rye{3}{The second compound below is the only instance in which the 音読み becomes voiced.}
\pronsrthreecone & --- & --- \\\hline
\end{prons}

\begin{remprons}
\hooglig{Reach} for the \comkun{italic} \comon{shield}.\\\hline
\end{remprons}

%{\middel}
% &  ()\\\hline
\begin{somme}
至る & \textit{reach; extend}  & いた{\middel}る \mbox{(ita{\middel}ru)}\\\hline
冬至 & winter + reach = \textit{winter solstice}  & トウ{\middel}ジ \mbox{(T\={O}{\middel}JI)}\\\hline
\notyet{至る所}  & reach + location = \textit{everywhere}  & いた{\middel}る　ところ \mbox{(ita{\middel}ru tokoro)}\\\hline
\notyet{至急} & reach + hasten = \textit{urgent}  & シ{\middel}キュウ \mbox{(SHI{\middel}KY\={U})}\\\hline
\end{somme}

\begin{decomp}{6}
この&道&は&公園&に&至る。\\\hline
この&みち&は&コウエン&に&いたる\\\hline
this&road&は&public park&に&arrive\\\hline
\multicolumn{6}{|r|}{\textit{This road goes to the park.}}\\\hline
\end{decomp}
\end{multicols}

\begin{decomp}{6}
英語&は&世界中&至る所&で&使われている。\\\hline
エイゴ&は&セカイジュウ&いたる　ところ&で&つかわれている\\\hline
English&は&throughout the world&everywhere&で&to use\\\hline
\multicolumn{6}{|r|}{\textit{English is used in every part of the world.}}\\\hline
\end{decomp}

\pagebreak
%%--------------------------------------------------------------------
%14-------------------------------------------------------------------

\begin{multicols}{2}
\begin{glyph}{Compare (比)}{compare}\label{glyph:compare}
Compare.\\\hline
The sense of \textit{ratio} is also present in a few compounds.
\end{glyph}
\gindex{Compare}{比}\strindex{4}{比}

\begin{comps}
\comprtwocone & 比 \\\hline
\comprthreecone & Compare/Side-by-side \\\hline
\end{comps}

\begin{remglyph}
Part of the 漢字 radicals (\#{81}).\\\hline
\end{remglyph}

\begin{prons}
\pronsrtwocone & \textbf{ヒ (HI)} & \textbf{くら (kura)}\\\hline
\pronsrthreecone & --- & ---\\\hline
\end{prons}

\begin{remprons}
\hooglig{Compare} the \comon{healing} \comkun{curators}.\\\hline
\end{remprons}

%{\middel}
% &  ()\\\hline
\begin{somme}
比べる& \textit{compare}  & くら{\middel}べる \mbox{(kura{\middel}beru)}  \\\hline
背比べ& back + compare = \textit{compare heights}  & せい　くら{\middel}べ \mbox{(sei kura{\middel}be)}  \\\hline
見比べる& see + compare = \textit{comparison}  & み　くら{\middel}べる \mbox{(mi kura{\middel}beru)}  \\\hline
\notyet{対比}& versus + compare = \textit{contrast; compare}  & タイ{\middel}ヒ \mbox{(TAI{\middel}HI)}  \\\hline
\end{somme}
\end{multicols}

\begin{decomp}{6}
僕&と&兄&を&比べないで&下さい。\\\hline
ボク&と&あに&を&くらべないで&ください\\\hline
me&と&brother&を&don't compare&please\\\hline
\multicolumn{6}{|r|}{\textit{Please don't compare me with my brother.}}\\\hline
\end{decomp}

%%--------------------------------------------------------------------
%15-------------------------------------------------------------------
\begin{multicols}{2}
\begin{glyph}{Measure (計)}{measure}\label{glyph:measure}
Measure.  Estimate.
\end{glyph}
\gindex{Measure}{計}\strindex{9}{計}

\begin{comps}
\comprtwocone & 言 + 十 \\\hline
\comprthreecone & Say + Ten \\\hline
\end{comps}

\begin{remglyph}
\hooglig{Measure} up to the \remcom{speech} of the \remcom{scarecrow}.\\\hline
\end{remglyph}

\begin{prons}
\pronsrtwocone & \textbf{ケイ (KEI)} & \textbf{はか (haka)}\\\hline
\pronsrthreecone & --- & --- \\\hline
\end{prons}

\begin{remprons}
\hooglig{Measure} up to the \comkun{HK's} \comon{catering} requests.\\\hline
\end{remprons}

%{\middel}
% &  ()\\\hline
\begin{somme}
\rye{3}{Don't forget that the reading of 時 in the third compound is irregular.}
計る  & \textit{measure; estimate}  & はか{\middel}る \mbox{(haka{\middel}ru)}\\\hline
会計  & meet + measure = \textit{accounts; accounting}  & カイ{\middel}ケイ \mbox{(KAI{\middel}KEI)}\\\hline
時計  & time + measure = \textit{watch; clock}  & と{\middel}ケイ \mbox{(to{\middel}KEI)}\\\hline
\notyet{温度計}  & warm + degree + measure = \textit{thermometer}  & オン{\middel}ド{\middel}ケイ \mbox{(ON{\middel}DO{\middel}KEI)}  \\\hline
\end{somme}

\begin{decomp}{6}
僕&の&時計&は&どこ&だ。\\\hline
ボク&の&とケイ&は&どこ&だ\\\hline
I&の&watch&は&where&be/is\\\hline
\multicolumn{6}{|r|}{\textit{Where is my watch?}}\\\hline
\end{decomp}

\begin{decomp}{3}
血圧&を&計りましょう。\\\hline
ケツアツ&を&はかりましょう\\\hline
blood pressure&を&to measure\\\hline
\multicolumn{3}{|r|}{\textit{Let me take your blood pressure.}}\\\hline
\end{decomp}
\end{multicols}

\pagebreak
%%--------------------------------------------------------------------
%16-------------------------------------------------------------------
\begin{multicols}{2}
\begin{glyph}{Location (所)}{location}\label{glyph:location}
Location.
\end{glyph}
\gindex{Location}{所}\strindex{8}{所}

\begin{comps}
\comprtwocone & 戸 + 斤\\\hline
\comprthreecone & Door + Axe\\\hline
\end{comps}

\begin{remglyph}
\hooglig{Locate} the \remcom{axed} \remcom{door}.\\\hline
\hooglig{Locate} the \remcom{door} with an \remcom{axe} on it.\\\hline
\end{remglyph}

\begin{prons}
\pronsrtwocone & \textbf{ショ (SHO)} & \textbf{ところ (tokoro)}\\\hline
\rye{3}{In the first position, ショ is by far the more common reading.\\\hline
In the second position, the 訓読み is always voiced, with the exceptions of only a few words in which this 漢字 is preceded by a verb in its る form (such as 至る所, the third example in Entry \ref{glyph:reach}, page \pageref{glyph:reach}).\\\hline
The 音読み, by contrast is almost never voiced in second position; the only time you are likely to encounter it voiced here is in the second example below.\\\hline
In the third position and beyond, however, there is no way to predict on first encountering a compound whether or not the 音読み is voiced (the 訓読み never appears after the second position).}
\pronsrthreecone & --- & --- \\\hline
\end{prons}

\begin{remprons}
The \hooglig{location} of the \comkun{tokoloshe} \comon{shop}.\\\hline
{\warn}In African folklore, a \textit{tokoloshe} is a mischievous and lascivious hairy water sprite.\\\hline
\end{remprons}

%{\middel}
% &  ()\\\hline
\begin{somme}
所 & \textit{location; place}  & ところ \mbox{(tokoro)}  \\\hline
近所 & near + location = \textit{neighbourhood; vicinity}  & キン{\middel}ジョ \mbox{(KIN{\middel}JO)}  \\\hline
\notyet{台所} & platform + location = \textit{kitchen}  & ダイ{\middel}どころ \mbox{(DAI{\middel}dokoro)}  \\\hline
\notyet{場所} & place + location = \textit{place}  & ば{\middel}ショ \mbox{(ba{\middel}SHO)}  \\\hline
\end{somme}

\begin{decomp}{5}
台所&に&猫&が&いる。\\\hline
ダイどころ&に&ねこ&が&いる\\\hline
kitchen&に&cat&が&to be\\\hline
\multicolumn{5}{|r|}{\textit{There is a cat in the kitchen.}}\\\hline
\end{decomp}

\begin{decomp}{4}
場所&を&教えて&下さい。\\\hline
ばショ&を&おしえて&ください\\\hline
place&を&to inform&please\\\hline
\multicolumn{4}{|r|}{\textit{Please tell me your location.}}\\\hline
\end{decomp}
\end{multicols}
%%--------------------------------------------------------------------
%17-------------------------------------------------------------------
\begin{multicols}{2}
\begin{glyph}{No. (第)}{number}\label{glyph:number}
This 漢字 is usually used as a prefix for numbers, much as English uses \textit{No.}, or \textit{\#}.\\\hline
In a few compounds it can also mean \textit{exam}.
\end{glyph}
\gindex{Number}{第}\strindex{11}{第}

\begin{comps}
\comprtwocone & 竹 + 弔 + ノ \\\hline
\comprthreecone & Bamboo + Mourn (\ref{glyph:mourn}) + NO/slash \\\hline
\end{comps}

\begin{remglyph}
\hooglig{Numerous} people \remcom{mourned} the \remcom{slashed} \remcom{bamboo}.\\\hline
\end{remglyph}

\begin{prons}
\pronsrtwocone & \textbf{ダイ (DAI)} & --- \\\hline
\pronsrthreecone & --- & --- \\\hline
\end{prons}

\begin{remprons}
\hooglig{Numerous} \comon{dyes}\ldots\\\hline
\end{remprons}

%{\middel}
% &  ()\\\hline
\begin{somme}
\rye{3}{Recall that `ease' + `complete' in the second compound below gives us \textit{safety}.}
第一  & number + one = \textit{No. 1}  & ダイ{\middel}イチ \mbox{(DAI{\middel}ICHI)}\\\hline
安全第一  & safety + number + one = \textit{``safety first''}  & アン{\middel}ゼン{\middel}ダイ{\middel}イチ \mbox{(AN{\middel}ZEN{\middel}DAI{\middel}ICHI)}\\\hline
\notyet{次第}  & next + Number = \textit{order; circumstances}  & シ{\middel}ダイ \mbox{(SHI{\middel}DAI)}\\\hline
\end{somme}

\begin{decomp}{8}
日本&は&天下&第一&の&国&が&ある。\\\hline
ニホン&は&テンカ&ダイイチ&の&くに&が&ある\\\hline
Japan&は&whole world&best&の&country&が&ある\\\hline
\multicolumn{8}{|r|}{\textit{Japan is the best country under the sun.}}\\\hline
\end{decomp}
\end{multicols}
\pagebreak
%%--------------------------------------------------------------------
%18-------------------------------------------------------------------

\begin{multicols}{2}
\begin{glyph}{Gist (旨)}{gist}\label{glyph:gist}
Gist.\\\hline
Though rarely encountered on its own, this 漢字 figures in other more common characters.
\end{glyph}
\gindex{Gist}{旨}\strindex{6}{旨}

\begin{comps}
\comprtwocone & ヒ + 日 \\\hline
\comprthreecone & Spoon/Ladle + Sun \\\hline
\end{comps}

\begin{remglyph}
The \hooglig{gist} of the \remcom{spoon} is as clear as \remcom{day}.\\\hline
\end{remglyph}

\begin{prons}
\pronsrtwocone & --- & \textbf{むね (mune)}\\\hline
\pronsrthreecone & シ (SHI) & --- \\\hline
\end{prons}

\begin{remprons}
The \hooglig{gist} is that the \comkun{moon ebbs} away to \ucomon{shield} itself from the sun.\\\hline
\end{remprons}

%{\middel}
% &  ()\\\hline
\begin{somme}
旨&\textit{gist; meaning}&むね \mbox{(mune)}  \\\hline
\end{somme}

\begin{decomp}{5}
あなた&の&主旨&は&分かる。\\\hline
あなた&の&シュシ&は&わかる\\\hline
you&の&point&は&understand\\\hline
\multicolumn{5}{|r|}{\textit{I see your point.}}\\\hline
\end{decomp}
\end{multicols}

%%--------------------------------------------------------------------
%19-------------------------------------------------------------------
\begin{multicols}{2}
\begin{glyph}{Wear (着)}{wear}\label{glyph:wear}
This important 漢字 has two main meanings:  \textit{to wear}, and \textit{to arrive}.\\\hline
It may help to think that, in a figurative sense, a person \textit{wears} their surroundings when they arrive at a place.\\\hline
This 漢字, in fact, consistently has this idea of \textit{things coming together} in certain ways.
\end{glyph}
\gindex{Wear}{着}\strindex{12}{着}

\begin{comps}
\comprtwocone & 羊 + 目 + 丿 \\\hline
\comprthreecone & Sheep + Eye + Slash/NO\\\hline
\end{comps}

\begin{remglyph}
\hooglig{Arrive} at the \remcom{grass} \remcom{king} and \remcom{slash} his \remcom{eye}.\\\hline
\hooglig{Arrive} with \remcom{sheep} \remcom{eyes} that are \remcom{not} blemished.\\\hline
\end{remglyph}

\begin{prons}
\pronsrtwocone & \textbf{チャク (CHAKU)} & \textbf{き、つ (ki, tsu)}\\\hline
\rye{3}{As the famous word in the third example makes clear, き is used when this 漢字 refers to \textit{wearing} clothes; it is always voiced in the second and third position.\\\hline
つ is used in the sense of \textit{arrive}, and is never voiced in the second position.\\\hline
チャク is the most common of these readings and has the overarching meaning of \textit{things coming together}.}
\pronsrthreecone & ジャク (JAKU) & --- \\\hline
\end{prons}

\begin{remprons}
\hooglig{Weary} \comon{Chuck} \comkun{keeps} \ucomon{jucking} the \comkun{tsunami} away.\\\hline
\end{remprons}

%{\middel}
% &  ()\\\hline
\begin{somme}
着る  & \textit{wear}  & き{\middel}る \mbox{(ki{\middel}ru)}  \\\hline
着く  & \textit{arrive}  & つ{\middel}く \mbox{(tsu{\middel}ku)}  \\\hline
着物  & wear + thing = \textit{kimono}  & き{\middel}もの \mbox{(ki{\middel}mono)}  \\\hline
\notyet{接着}  & join + come together = \textit{adhesion}  & セッ{\middel}チャク \mbox{(SET{\middel}CHAKU)}  \\\hline
\end{somme}

\begin{decomp}{5}
私&は&駅&に&着いた。\\\hline
わたし&は&エキ&に&ついた\\\hline
I&は&station&に&arrive\\\hline
\multicolumn{5}{|r|}{\textit{I arrived at the station.}}\\\hline
\end{decomp}

\begin{decomp}{5}
何&と&素晴らしい&着物&でしょう。\\\hline
なに&と&すばらしい&きもの&でしょう\\\hline
what&と&wonderful&dress&でしょう\\\hline
\multicolumn{5}{|r|}{\textit{What a heavenly dress!}}\\\hline
\end{decomp}
\end{multicols}

\pagebreak
%%--------------------------------------------------------------------
%20-------------------------------------------------------------------

\begin{multicols}{2}
\begin{glyph}{Lack (欠)}{lack}\label{glyph:lack}
Lack.
\end{glyph}
\gindex{Lack}{欠}\strindex{4}{欠}

\begin{comps}
\comprtwocone & 欠 \\\hline
\comprthreecone & Gap/Lack/Yawn \\\hline
\end{comps}

\begin{remglyph}
Part of the 漢字 radicals (\#{76}).\\\hline
\end{remglyph}

\begin{prons}
\pronsrtwocone & \textbf{ケツ (KETSU)} & \textbf{か (ka)}\\\hline
\pronsrthreecone & --- & --- \\\hline
\end{prons}

\begin{remprons}
A \hooglig{lack} of \comon{kets} \comkun{customs}.\\\hline
{\warn}\textit{Ket} people represent an ethnic group in Siberia.\\\hline
\end{remprons}

%{\middel}
% &  ()\\\hline
\begin{somme}
欠く& \textit{lack}  & か{\middel}く \mbox{(ka{\middel}ku)}\\\hline
\notyet{欠点}& lack + point = \textit{defect; shortcoming}  & ケッ{\middel}テン \mbox{(KET{\middel}TEN)}\\\hline
\notyet{欠員}& lack + member = \textit{vacancy; post}  & ケツイン \mbox{(KETSU{\middel}IN)}\\\hline
\end{somme}

\begin{decomp}{5}
欠点&なき&人&は&なし。\\\hline
ケッテン&なき&ひと&は&なし\\\hline
fault&non existent&person&は&non existent\\\hline
\multicolumn{5}{|r|}{\textit{No man is without his faults.}}\\\hline
\end{decomp}

\begin{decomp}{6}
他人&の&欠点&を&探す&な。\\\hline
タニン&の&ケッテン&を&さがす&な\\\hline
others&の&fault&を&to find&な\\\hline
\multicolumn{6}{|r|}{\textit{Don't find fault with others.}}\\\hline
\end{decomp}
\end{multicols}

%%--------------------------------------------------------------------
%21-------------------------------------------------------------------
\begin{multicols}{2}
\begin{glyph}{Show (示)}{show}\label{glyph:show}
To show.
\end{glyph}
\gindex{Show}{示}\strindex{5}{示}

\begin{comps}
\comprtwocone & 示 \\\hline
\comprthreecone & Altar/Display \\\hline
\end{comps}

\begin{remglyph}
Part of the 漢字 radicals (\#{113}).\\\hline
\end{remglyph}

\begin{prons}
\pronsrtwocone & \textbf{ジ (JI)} & \textbf{しめ (shime)}\\\hline
\pronsrthreecone & シ (SHI)  & ---\\\hline
\end{prons}

\begin{remprons}
The \hooglig{show} \comkun{she met} was \comon{jigged} with \ucomon{shields}.\\\hline
\end{remprons}

%{\middel}
% &  ()\\\hline
\begin{somme}
示す& \textit{show}  & しめ{\middel}す \mbox{(shime{\middel}su)}\\\hline
暗示& dark + show = \textit{hint}  & アン{\middel}ジ \mbox{(AN{\middel}JI)}\\\hline
\notyet{指示}& finger + show = \textit{indication; direction}  & シ{\middel}ジ \mbox{(SHI{\middel}JI)}\\\hline
\notyet{展示}& display + show = \textit{display; exhibition}  & テン{\middel}ジ \mbox{(TEN{\middel}JI)}\\\hline
\end{somme}

\begin{decomp}{7}
雪&は&冬&の&到来&を&示す。\\\hline
ゆき&は&ふゆ&の&トウライ&を&しめす\\\hline
snow&は&winter&の&coming&を&show\\\hline
\multicolumn{7}{|r|}{\textit{Snow indicates the coming of winter.}}\\\hline
\end{decomp}
\end{multicols}

\begin{decomp}{9}
彼&の&言葉&は&何&を&暗示している&の&か。\\\hline
かれ&の&ことば&は&なに&を&アンジしている&の&か\\\hline
he&の&word&は&what&を&hint&の&か\\\hline
\multicolumn{9}{|r|}{\textit{What do his words imply?}}\\\hline
\end{decomp}

\pagebreak
%%--------------------------------------------------------------------
%22-------------------------------------------------------------------

\begin{multicols}{2}
\begin{glyph}{Serve (仕)}{serve}\label{glyph:serve}
Serve.  Work.
\end{glyph}
\gindex{Serve}{仕}\strindex{5}{仕}

\begin{comps}
\comprtwocone & 亻 + 士 \\\hline
\comprthreecone & Person + Gentleman \\\hline
\end{comps}

\begin{remglyph}
\hooglig{Serve} that \remcom{person} as if he is a \remcom{gentlemen}.\\\hline
\end{remglyph}

\begin{prons}
\pronsrtwocone & \textbf{シ (SHI)} & --- \\\hline
\pronsrthreecone & ジ (JI) & つか (tsuka) \\\hline
\end{prons}

\begin{remprons}
I \hooglig{serve} by \ucomon{jigging} \comon{shields} and \ucomkun{tsukas}.\\\hline
{\warn}The \textit{tsuka} is the hilt or handle of a Japanese sword.\\\hline
\end{remprons}

%{\middel}
% &  ()\\\hline
\begin{somme}
\rye{3}{{\diff}This 漢字 is unusual in that its 音読み almost always forms compounds with the 訓読み of other characters, as in the examples below.}
仕方 & serve + direction = \textit{way; method}  & シ{\middel}かた \mbox{(SHI{\middel}kata)}\\\hline
仕事 & serve + matter = \textit{job; work}  & シ{\middel}ごと \mbox{(SHI{\middel}goto)}\\\hline
仕切り & serve + cut = \textit{partition; dividing line}  & シ{\middel}きり \mbox{(SHI{\middel}kiri)}\\\hline
仕上げる & serve + upper = \textit{finish; complete}  & シ あ{\middel}げる \mbox{(SHI a{\middel}geru)}\\\hline
\end{somme}

\begin{decomp}{4}
仕方&が&無い&よ。\\\hline
シかた&が&ない&よ\\\hline
way&が&non-existent&よ\\\hline
\multicolumn{4}{|r|}{\textit{There's no choice.}}\\\hline
\end{decomp}

\begin{decomp}{6}
彼&は&仕事&の&鬼&だ。\\\hline
かれ&は&シごと&の&おに&だ\\\hline
he&は&work&の&demon&be/is\\\hline
\multicolumn{6}{|r|}{\textit{He is an eager beaver.}}\\\hline
\end{decomp}
\end{multicols}

%%--------------------------------------------------------------------
%23-------------------------------------------------------------------
\begin{multicols}{2}
\begin{glyph}{Heavy (重)}{heavy}\label{glyph:heavy}
Heavy, both in the literal and figurative sense.
\end{glyph}
\gindex{Heavy}{重}\strindex{9}{重}

\begin{comps}
\comprtwocone &  丿 + 里  \\\hline
\comprthreecone & Slash/NO + Hamlet \ref{glyph:hamlet} \\\hline
\end{comps}

\begin{remglyph}
\hooglig{Heavily} \remcom{slash} the \remcom{hamlet}.\\\hline
\end{remglyph}

\begin{prons}
\pronsrtwocone & \textbf{ジュウ (J\={U})} & \textbf{おも (omo)}\\\hline
\rye{3}{ジュウ is by far the more common reading.}
\pronsrthreecone & チョウ (CH\={O}) & かさ、え (kasa, e) \\\hline
\end{prons}

\begin{remprons}
\hooglig{Heavy} \ucomkun{edging} \ucomon{chores} for \ucomkun{casanova} due to his \comon{juice} stains that required \comkun{omo}.\\\hline
\end{remprons}

%{\middel}
% &  ()\\\hline
\begin{somme}
重い  & \textit{heavy}  & おも{\middel}い \mbox{(omo{\middel}i)}  \\\hline
重大  & heavy + large = \textit{important; grave}  & ジュウ{\middel}ダイ \mbox{(J\={U}{\middel}DAI)}  \\\hline
重要  & heavy + essential = \textit{important; essential}  & ジュウ{\middel}ヨウ \mbox{(J\={U}{\middel}Y\={O})}  \\\hline
\end{somme}

\begin{decomp}{5}
彼&は&口&が&重い。\\\hline
かれ&は&くち&が&おもい\\\hline
he&は&mouth&が&heavy\\\hline
\multicolumn{5}{|r|}{\textit{He is a man of few words.}}\\\hline
\end{decomp}

\begin{decomp}{6}
結婚&は&重大&な&問題&だ。\\\hline
ケッコン&は&ジュウダイ&な&モンダイ&だ\\\hline
marriage&は&serious&な&matter&be/is\\\hline
\multicolumn{6}{|r|}{\textit{Marriage is a serious matter.}}\\\hline
\end{decomp}
\end{multicols}

\pagebreak

%%--------------------------------------------------------------------
%24-------------------------------------------------------------------

\begin{multicols}{2}
\begin{glyph}{Light (軽)}{light}\label{glyph:light}
Light, in terms of weight or feeling.
\end{glyph}
\gindex{Light}{軽}\strindex{12}{軽}

\begin{comps}
\comprtwocone&車 + 又 + 土\\\hline
\comprthreecone & Car + Again + Earth\\\hline
\end{comps}

\begin{remglyph}
The \hooglig{light} \remcom{car} is \remcom{again} on \remcom{soil}.\\\hline
\end{remglyph}

\begin{prons}
\pronsrtwocone&\textbf{ケイ (KEI)}&\textbf{かる (karu)}\\\hline
\rye{3}{The 訓読み is voiced when not in the first position, as seen in
	the second and third compounds (the lattter a mix of 音読み and 訓読み).}
\pronsrthreecone&---&かろ (karo)\\\hline
\end{prons}

\begin{remprons}
\hooglig{Light} \comkun{caruncles} were \comon{catering} for the
\ucomkun{carousel}.\\\hline
{\warn}\textit{Caruncle} is a part of the human eye anatomy.\\\hline{\warn}\textit{Carousel} is a merry-go-round.\\\hline
\end{remprons}

%{\middel}
% &  ()\\\hline
\begin{somme}
軽い&\textit{light}&かる{\middel}い \mbox{(karu{\middel}i)}\\\hline
手軽&hand + light = \textit{simple; easy}&て{\middel}がる \mbox{(te{\middel}garu)}\\\hline
気軽&sprit + light = \textit{lighthearted}&キ{\middel}がる \mbox{(KI{\middel}garu)}\\\hline
\notyet{軽食}&light + eat = \textit{light meal; snack}& ケイ{\middel}ショク \mbox{(KEI{\middel}SHOKU)}\\\hline
\end{somme}

\begin{decomp}{7}
軽食&の&出来る&カフェ&が&あります&か。\\\hline
ケイショク&の&できる&カフェ&が&あります&か\\\hline
light snack&の&be able to&cafe&が&to be&か\\\hline
\multicolumn{7}{|r|}{\textit{Is there a cafe where I can have a light meal?}}\\\hline
\end{decomp}
\end{multicols}

%%--------------------------------------------------------------------
%25-------------------------------------------------------------------

\begin{multicols}{2}
\begin{glyph}{Match (合)}{match}\label{glyph:match}
Match.  Combine.\\\hline
This is a 漢字 you will encounter frequently, as it forms many compound verbs (such as in the third example) that suggest some kind of \textit{reciprocal action} between people or things.  These words are often translated into English with the phrase ``each other'': 押し合う (お{\middel}し　あ{\middel}う) \textit{to push each other}, is one example.\\\hline
This 漢字 can also act like 買 in absorbing its ひらがな accompaniment in certain compounds.
\end{glyph}
\gindex{Match}{合}\strindex{6}{合}

\begin{comps}
\comprtwocone & (亼 = 𠆢 + 一) + 口 \\\hline
\comprthreecone & (Assembly = Person + One) + Mouth \\\hline
\end{comps}

\begin{remglyph}
\hooglig{Match} that \remcom{assembly's} \remcom{vampire} speech.\\\hline
\end{remglyph}

\begin{prons}
\pronsrtwocone & \textbf{ゴウ (G\={O})} & \textbf{あ (a)}\\\hline
\rye{3}{ゴウ is the most common choice in the first position, but both it and あ occur frequently in second position.}
\pronsrthreecone & ガッ、カッ (GA-, KA-) & --- \\\hline
\rye{3}{These 2 readings only occur in a few words, and always double up with the 音読み of the 漢字 that follows them (as the small カタカナ ッ indicates).\\\hline
The compound 合宿 (ガッ{\middel}シュク) \textit{to lodge together}, from the 漢字 `match' and `lodging' is an example.}
\end{prons}

\begin{remprons}
The \hooglig{match} \comkun{update} was filled with \ucomon{cuss-worthy} \ucomon{gushing} \comon{gore}.\\\hline
\end{remprons}

%{\middel}
% &  ()\\\hline
\begin{somme}
合う (intr.)  & \textit{match; fit}  & あ{\middel}う \mbox{(a{\middel}u)}  \\\hline
合わせる (tr.)  & \textit{put together; combine}  & あ{\middel}わせる \mbox{(a{\middel}waseru)}  \\\hline
押し合う (tr.) & push + match = \textit{to push each other} & お{\middel}し　あ{\middel}う \mbox{(o{\middel}shi a{\middel}u)} \\\hline
話し合う  & speak + match = \textit{discuss; talk}  & はな{\middel}し　あ{\middel}う \mbox{(hana{\middel}shi a{\middel}u)}  \\\hline
合意  & match + mind = \textit{mutual agreement}  & ゴウ{\middel}イ \mbox{(G\={O}{\middel}I)}  \\\hline
合宿 & match + lodging = \textit{to lodge together} & ガッ{\middel}シュク \mbox{(GAS{\middel}SHUKU)}\\\hline
\end{somme}

\begin{decomp}{3}
合言葉&を&言う。\\\hline
あいことば&を&いう\\\hline
password&を&say\\\hline
\multicolumn{3}{|r|}{\textit{Give the password.}}\\\hline
\end{decomp}

\begin{decomp}{4}
会合&は&先月&あった。\\\hline
カイゴウ&は&センゲツ&あった\\\hline
meeting&は&last month&happened\\\hline
\multicolumn{4}{|r|}{\textit{The meeting was last month.}}\\\hline
\end{decomp}
\end{multicols}

\pagebreak
%%--------------------------------------------------------------------
